% figures/w02-tuning-order.svg — PID 튜닝 순서. 한 번에 하나씩.
%
%   bash _tools/tikz2svg.sh figures/src/w02-tuning-order.tex figures/w02-tuning-order.svg
%
% 2주차 §2-12 의 그림. 사용자 요청(2026-09-21): "PID 튜닝 하는 순서도 정리해줘봐....
% 최대한 쉽게". 순서의 근거는 셋이다.
%   - 제어조교(Ctrl튜브) "PID 제어기 짬튜닝" — P 를 먼저, 단위로 크기를 정하고, P 응답으로
%     시스템을 읽는다. I 는 남는 오차가 있을 때만. 목표값은 부드럽게.
%   - MATLAB Tech Talk "Understanding PID Control" 4편·6편 — 먼저 요구조건과 플랜트가
%     "얌전한가" 를 본다. 끝에는 언제나 손으로 다듬는다. 제어 입력을 반드시 본다.
%   - 캡스톤디자인 6주차 F-6 — P → D → I → 안티와인드업.
% 오른쪽 열의 수는 W02_I_tuning_by_hand 가 이 플랜트에서 실제로 잰 값이다.
%
% 폭 약 176 mm.
\documentclass[border=5pt]{standalone}
% gnc-style.tex — 이 볼트의 TikZ 그림이 공유하는 스타일.
%
% 저널 그림 규약을 스타일 하나로 굳혀 둔다. 그림마다 선 굵기나 화살촉을
% 다시 정하지 않는다 — 그것이 그림이 제각각으로 보이는 원인이다.
%
% 쓰는 법:  % gnc-style.tex — 이 볼트의 TikZ 그림이 공유하는 스타일.
%
% 저널 그림 규약을 스타일 하나로 굳혀 둔다. 그림마다 선 굵기나 화살촉을
% 다시 정하지 않는다 — 그것이 그림이 제각각으로 보이는 원인이다.
%
% 쓰는 법:  % gnc-style.tex — 이 볼트의 TikZ 그림이 공유하는 스타일.
%
% 저널 그림 규약을 스타일 하나로 굳혀 둔다. 그림마다 선 굵기나 화살촉을
% 다시 정하지 않는다 — 그것이 그림이 제각각으로 보이는 원인이다.
%
% 쓰는 법:  \input{gnc-style}  (figures/src/ 안에서)

\usepackage{tikz}
\usepackage{amsmath}
\usepackage{xcolor}
\usetikzlibrary{arrows.meta, positioning, calc, fit, backgrounds, decorations.pathmorphing}

% ---- 색 --------------------------------------------------------------------
% 색은 강조 하나에만 쓴다. 나머지는 검정.
\definecolor{ink}{HTML}{1A1A1A}
\definecolor{accent}{HTML}{7C3AED}
\definecolor{warn}{HTML}{B91C1C}
\definecolor{soft}{HTML}{666666}

% ---- 선과 화살촉 -----------------------------------------------------------
\tikzset{
  % 신호선. 화살촉은 작고 뾰족한 것 하나뿐이다.
  sig/.style   = {ink, line width=0.5pt, -{Stealth[length=4pt, width=3pt]}},
  wire/.style  = {ink, line width=0.5pt},
  % 강조 경로 — 파선으로 구분한다. 흑백 인쇄에서도 살아야 하기 때문이다.
  asig/.style  = {accent, line width=0.5pt, densely dashed,
                  -{Stealth[length=4pt, width=3pt]}},
  awire/.style = {accent, line width=0.5pt, densely dashed},
  % 블록. 흰 바탕, 직각, 검은 테두리.
  blk/.style   = {draw=ink, line width=0.55pt, fill=white,
                  minimum width=17mm, minimum height=8mm, inner sep=2pt, font=\small},
  % 합산점
  sum/.style   = {draw=ink, line width=0.55pt, fill=white, circle,
                  inner sep=0pt, minimum size=4.6mm, font=\scriptsize},
  asum/.style  = {draw=accent, line width=0.55pt, fill=white, circle,
                  inner sep=0pt, minimum size=4.6mm, font=\scriptsize},
  % 분기점
  dot/.style   = {ink, fill=ink, circle, inner sep=0pt, minimum size=1.5mm},
  % 신호 이름 — 선 위에 작게
  sname/.style = {font=\small, inner sep=1.5pt},
  % 그림 안의 짧은 설명. 그림 안의 글은 최소로 둔다
  note/.style  = {font=\scriptsize, soft, inner sep=1.5pt},
  anote/.style = {font=\scriptsize, accent, inner sep=1.5pt},
  % 축
  axis/.style  = {ink, line width=0.5pt},
  tick/.style  = {font=\scriptsize, soft},
}

% 캡션 한 줄 — 그림 아래 가는 선과 함께
\newcommand{\gnccaption}[2]{%
  \node[anchor=north west, text width=#1, align=left, font=\scriptsize, ink]
        at (current bounding box.south west) {\vspace{2pt}#2};}
  (figures/src/ 안에서)

\usepackage{tikz}
\usepackage{amsmath}
\usepackage{xcolor}
\usetikzlibrary{arrows.meta, positioning, calc, fit, backgrounds, decorations.pathmorphing}

% ---- 색 --------------------------------------------------------------------
% 색은 강조 하나에만 쓴다. 나머지는 검정.
\definecolor{ink}{HTML}{1A1A1A}
\definecolor{accent}{HTML}{7C3AED}
\definecolor{warn}{HTML}{B91C1C}
\definecolor{soft}{HTML}{666666}

% ---- 선과 화살촉 -----------------------------------------------------------
\tikzset{
  % 신호선. 화살촉은 작고 뾰족한 것 하나뿐이다.
  sig/.style   = {ink, line width=0.5pt, -{Stealth[length=4pt, width=3pt]}},
  wire/.style  = {ink, line width=0.5pt},
  % 강조 경로 — 파선으로 구분한다. 흑백 인쇄에서도 살아야 하기 때문이다.
  asig/.style  = {accent, line width=0.5pt, densely dashed,
                  -{Stealth[length=4pt, width=3pt]}},
  awire/.style = {accent, line width=0.5pt, densely dashed},
  % 블록. 흰 바탕, 직각, 검은 테두리.
  blk/.style   = {draw=ink, line width=0.55pt, fill=white,
                  minimum width=17mm, minimum height=8mm, inner sep=2pt, font=\small},
  % 합산점
  sum/.style   = {draw=ink, line width=0.55pt, fill=white, circle,
                  inner sep=0pt, minimum size=4.6mm, font=\scriptsize},
  asum/.style  = {draw=accent, line width=0.55pt, fill=white, circle,
                  inner sep=0pt, minimum size=4.6mm, font=\scriptsize},
  % 분기점
  dot/.style   = {ink, fill=ink, circle, inner sep=0pt, minimum size=1.5mm},
  % 신호 이름 — 선 위에 작게
  sname/.style = {font=\small, inner sep=1.5pt},
  % 그림 안의 짧은 설명. 그림 안의 글은 최소로 둔다
  note/.style  = {font=\scriptsize, soft, inner sep=1.5pt},
  anote/.style = {font=\scriptsize, accent, inner sep=1.5pt},
  % 축
  axis/.style  = {ink, line width=0.5pt},
  tick/.style  = {font=\scriptsize, soft},
}

% 캡션 한 줄 — 그림 아래 가는 선과 함께
\newcommand{\gnccaption}[2]{%
  \node[anchor=north west, text width=#1, align=left, font=\scriptsize, ink]
        at (current bounding box.south west) {\vspace{2pt}#2};}
  (figures/src/ 안에서)

\usepackage{tikz}
\usepackage{amsmath}
\usepackage{xcolor}
\usetikzlibrary{arrows.meta, positioning, calc, fit, backgrounds, decorations.pathmorphing}

% ---- 색 --------------------------------------------------------------------
% 색은 강조 하나에만 쓴다. 나머지는 검정.
\definecolor{ink}{HTML}{1A1A1A}
\definecolor{accent}{HTML}{7C3AED}
\definecolor{warn}{HTML}{B91C1C}
\definecolor{soft}{HTML}{666666}

% ---- 선과 화살촉 -----------------------------------------------------------
\tikzset{
  % 신호선. 화살촉은 작고 뾰족한 것 하나뿐이다.
  sig/.style   = {ink, line width=0.5pt, -{Stealth[length=4pt, width=3pt]}},
  wire/.style  = {ink, line width=0.5pt},
  % 강조 경로 — 파선으로 구분한다. 흑백 인쇄에서도 살아야 하기 때문이다.
  asig/.style  = {accent, line width=0.5pt, densely dashed,
                  -{Stealth[length=4pt, width=3pt]}},
  awire/.style = {accent, line width=0.5pt, densely dashed},
  % 블록. 흰 바탕, 직각, 검은 테두리.
  blk/.style   = {draw=ink, line width=0.55pt, fill=white,
                  minimum width=17mm, minimum height=8mm, inner sep=2pt, font=\small},
  % 합산점
  sum/.style   = {draw=ink, line width=0.55pt, fill=white, circle,
                  inner sep=0pt, minimum size=4.6mm, font=\scriptsize},
  asum/.style  = {draw=accent, line width=0.55pt, fill=white, circle,
                  inner sep=0pt, minimum size=4.6mm, font=\scriptsize},
  % 분기점
  dot/.style   = {ink, fill=ink, circle, inner sep=0pt, minimum size=1.5mm},
  % 신호 이름 — 선 위에 작게
  sname/.style = {font=\small, inner sep=1.5pt},
  % 그림 안의 짧은 설명. 그림 안의 글은 최소로 둔다
  note/.style  = {font=\scriptsize, soft, inner sep=1.5pt},
  anote/.style = {font=\scriptsize, accent, inner sep=1.5pt},
  % 축
  axis/.style  = {ink, line width=0.5pt},
  tick/.style  = {font=\scriptsize, soft},
}

% 캡션 한 줄 — 그림 아래 가는 선과 함께
\newcommand{\gnccaption}[2]{%
  \node[anchor=north west, text width=#1, align=left, font=\scriptsize, ink]
        at (current bounding box.south west) {\vspace{2pt}#2};}

\usetikzlibrary{shapes.geometric}
\tikzset{
  tstep/.style = {draw=ink, line width=0.55pt, fill=white, rounded corners=1mm,
                 text width=66mm, minimum height=10mm, inner sep=2.2mm,
                 font=\small, align=left},
  ask/.style  = {draw=accent, line width=0.6pt, fill=white, diamond, aspect=2.6,
                 inner sep=0.6mm, font=\small, align=center, text=accent},
  num/.style  = {font=\scriptsize, soft, text width=52mm, align=left, anchor=west},
  lbl/.style  = {font=\scriptsize, accent, inner sep=1pt},
}

\begin{document}
\begin{tikzpicture}[x=1mm, y=1mm]

\node[font=\small\bfseries, anchor=south] at (10,14)
     {Tuning a PID by hand: one gain at a time, each chosen from the last measurement};

% ---- 0 ---------------------------------------------------------------------
\node[tstep] (s0) at (0,0) {\textbf{0\quad Before any gain.} Write down what ``good'' means
  (overshoot, settling time, force available). Check the plant is well behaved:
  stable, nearly linear, little delay.};

% ---- 1 ---------------------------------------------------------------------
\node[tstep, below=6mm of s0] (s1) {\textbf{1\quad P only.} $K_i = K_d = 0$. Choose the size of
  $K_p$ from the units: a few times the stiffness the plant already has.};
\node[num] at ([xshift=6mm]s1.east) {$k = 2$ N/m, so $K_p = 5k = 10$ N/m};

% ---- 2 ---------------------------------------------------------------------
\node[tstep, below=6mm of s1] (s2) {\textbf{2\quad Read the P response.} The overshoot
  $M_p$ gives the damping ratio:\\[1mm]
  \hspace*{8mm}$\zeta = -\ln M_p \big/ \sqrt{\pi^2 + \ln^2 M_p}$};
\node[num] at ([xshift=6mm]s2.east) {$M_p = 38.8\,\%$, so $\zeta = 0.289$};

\node[ask, below=6mm of s2] (q2) {does it ring?};
\draw[sig] (s0) -- (s1);  \draw[sig] (s1) -- (s2);  \draw[sig] (s2) -- (q2);

% no: raise Kp
\node[tstep, text width=36mm, anchor=east] (up) at ([xshift=-12mm]q2.west)
     {\textbf{no}: slow and smooth. Raise $K_p$ and read again.};
\draw[asig] (q2.west) -- node[lbl, above] {no} (up.east);
\draw[asig] (up.north) |- ([yshift=-2mm]s1.west);

% ---- 3 ---------------------------------------------------------------------
\node[tstep, below=6mm of q2] (s3) {\textbf{3\quad Add D.} Raise $K_d$ step by step while the
  overshoot keeps falling. Stop where it stops falling. If the force starts to
  chatter, lower $N_f$.};
\draw[asig] (q2) -- node[lbl, right] {yes} (s3);
\node[num] at ([xshift=6mm]s3.east) {$K_d = 0,1,\dots,7$: $38.8 \to 8.60\,\%$ at $K_d = 6$, then up again};

\node[ask, below=6mm of s3] (q3) {error left?};
\draw[sig] (s3) -- (q3);

% ---- 4 ---------------------------------------------------------------------
\node[tstep, below=6mm of q3] (s4) {\textbf{4\quad Add I}, only because an error remains. Raise
  $K_i$ until the response is inside the band in time. Too much brings the overshoot back.};
\draw[asig] (q3) -- node[lbl, right] {yes} (s4);
\node[num] at ([xshift=6mm]s4.east) {error left $0.167$ m; $K_i = 8$: inside 1\,\% in $2.77$ s};
\node[tstep, text width=36mm, anchor=east] (noI) at ([xshift=-12mm]q3.west)
     {\textbf{no}: leave $K_i = 0$. An integral that is not needed only costs.};
\draw[asig] (q3.west) -- node[lbl, above] {no} (noI.east);

% ---- 5 ---------------------------------------------------------------------
\node[ask, below=6mm of s4] (q5) {force within\\the limit?};
\draw[sig] (s4) -- (q5);
%  I 가 필요 없으면 4 를 건너 4 와 5 사이의 연결선으로 들어간다.
%  With no error left, step 4 is skipped: join the connector between 4 and 5.
\coordinate (j45) at ($(s4.south)!0.5!(q5.north)$);
\draw[asig] (noI.south) |- (j45);
\node[tstep, below=6mm of q5] (s5) {\textbf{5\quad Look at the force, not only the response.}
  Smooth the setpoint to remove the derivative kick; switch on anti-windup ($K_b$);
  if still too much, lower the gains.};
\draw[asig] (q5) -- node[lbl, right] {no} (s5);
\node[num] at ([xshift=6mm]s5.east) {step: $129.6$ N against $30$ N available; smoothed: $13.9$ N};

% ---- 6 ---------------------------------------------------------------------
\node[tstep, below=6mm of s5] (s6) {\textbf{6\quad Finish by hand.} A PID Tuner or a
  Ziegler--Nichols rule gives a starting point, never the last word. Re-check
  step 5 after every change.};
\draw[sig] (s5) -- (s6);
%  힘이 한계 안이면 5 를 건너 6 으로 / within the limit: skip 5, go to 6
\coordinate (byp) at ([xshift=-6mm]s5.west |- q5.west);
\draw[asig] (q5.west) -- node[lbl, above] {yes} (byp) |- (s6.west);

\end{tikzpicture}
\end{document}
